% ---
% id: foundations/lem:exact-accumulation
% kind: lem
% title: "Exactness of a sum that rounds nowhere"
% summary: >-
%   A sequential sum rounds nowhere exactly when every exact partial sum is a
%   binary64 number, and then the computed sum equals the exact sum at every
%   step. A predicate stated on such a sum is a predicate on the declared values
%   themselves.
% prominence: main
% status: accepted
% uses: [def:declared-order, def:binary64]
% assumes: []
% refs: []
% provenance: newly-derived
% lean:
% ---

\begin{lemma}\label{lem:exact-accumulation}
Let $(x_1,\dots,x_n)$ be a finite sequence of elements of
$\mathbb{B}_{\mathbb{C}}$, or of matrices of one shape over
$\mathbb{B}_{\mathbb{C}}$, such as declared values or the computed inner sums
of a sum of sums, with sequential sum $(S_i)_{i=0}^{n}$ and computed
sequential sum $(\hat S_i)_{i=0}^{n}$ (Definition~\ref{def:declared-order}).
The following are equivalent:
\begin{enumerate}[label=(\alph*),leftmargin=2.4em]
\item the accumulation rounds nowhere;
\item $\hat S_i=S_i$ for every $0\le i\le n$;
\item for every $1\le i\le n$, the real and the imaginary part of every entry
  of $S_i$ are binary64 numbers.
\end{enumerate}
If they hold, the computed sum $\hat S_n$ is the exact sum
$\sum_{i=1}^{n}x_i$.
\end{lemma}

\begin{proof}
The map $\mathrm{fl}$ acts on the real and the imaginary part of every entry
separately (Definition~\ref{def:binary64}\,(ii)), and $S_i$ and $\hat S_i$ are
formed entry by entry, so it suffices to prove the lemma for a sequence of
real numbers $x_i\in\mathbb{B}$.

(a)$\Rightarrow$(b). By induction on $i$: $\hat S_0=0=S_0$, and if
$\hat S_{i-1}=S_{i-1}$, then (a) gives
$\hat S_i=\hat S_{i-1}+x_i=S_{i-1}+x_i=S_i$.

(b)$\Rightarrow$(c). Let $1\le i\le n$. The value
$\hat S_i=\mathrm{fl}(\hat S_{i-1}+x_i)$ lies in $\mathbb{B}\cup\{\pm\infty\}$,
by the first property of $\mathrm{fl}$ where $\hat S_{i-1}$ is finite and by
its definition on $\pm\infty$ otherwise. By (b) it equals the real number
$S_i$, so $S_i\in\mathbb{B}$.

(c)$\Rightarrow$(a). We show by induction on $i$ that $\hat S_i=S_i$ and that
the $i$-th addition is exact. For $i=0$, $\hat S_0=S_0=0$. Let $1\le i\le n$
and $\hat S_{i-1}=S_{i-1}$. Then $\hat S_{i-1}$ is finite, and
$\hat S_{i-1}+x_i=S_i\in\mathbb{B}$ by (c), so the second property of
$\mathrm{fl}$ gives $\hat S_i=\mathrm{fl}(S_i)=S_i=\hat S_{i-1}+x_i$: the
$i$-th addition is exact, and $\hat S_i=S_i$.

Finally, if (b) holds then $\hat S_n=S_n=\sum_{i=1}^{n}x_i$ by
\eqref{eq:seqsum}.
\end{proof}
